\subsection{The pencil, the fixed part, and the rank}

We work over an algebraically closed field $k$.  Let $X/k$ be a smooth,
connected projective surface, let $H$ be a very ample divisor, and take a
pencil with smooth geometrically connected generic member and a transverse
base locus $Z$ consisting of $n=H^2$ distinct points.  Write
\[
 \pi:Y=\operatorname{Bl}_Z(X)\longrightarrow X,\qquad
 f:Y\longrightarrow\mathbf P^1_k,\qquad K=k(t),\quad C=Y_K.
\]
The exceptional curves are sections; choose one and call it $O$.
Let $F$ be the fibre class.  Write $J=\operatorname{Pic}^0_{C/K}$,
$T=\operatorname{Pic}^0_{Y/k,\mathrm{red}}$, and let
$\tau:T_K\to J$ be restriction.  Its image is denoted by $B$, and
$q:J\to J'=J/B$ is the mobile quotient.  In this section
$\rho(J')$ means $\dim_{\mathbf Q}(J'(K)\otimes\mathbf Q)$, exactly the
Mordell--Weil rank intended on cote 130; it is not the Picard number of
the abelian variety $J'$.  In positive characteristic we always use the
reduced Picard variety, and we remove N\'eron--Severi torsion when discussing
lattices.

\begin{lemma}[Proved: the fixed part and rational-point comparison]
\label{D:fixed}
The map $\tau$ is the $K/k$-trace map, up to its possible finite kernel.
Restriction and quotient induce an isomorphism
\[
 \bigl(J(K)/\tau T(k)\bigr)\otimes\mathbf Q
       \xrightarrow{\sim}J'(K)\otimes\mathbf Q.
\]
\end{lemma}
\begin{proof}
We recall directly why the Picard variety of the surface supplies the
fixed part in this situation.  For a constant abelian variety $T_1/k$
and a homomorphism $g:T_{1,K}\to J$, let $a:C\to J$ be the Abel--Jacobi
map based at $O_K$ and let $\lambda:J\to J^\vee$ be the positive principal
polarization.  The composite
\[
 C\xrightarrow{a}J\xrightarrow{\lambda}J^\vee
     \xrightarrow{g^\vee}T_{1,K}^\vee
\]
gives a rational $k$-map $Y\dashrightarrow T_1^\vee$.
A rational map from a smooth variety to an abelian variety extends to
a morphism \cite[Theorem 3.1]{MilneAV}.  It vanishes along $O$, so, after
choosing a $k$-point on $O$ as base point, the Albanese universal property
gives $u:\operatorname{Alb}(Y)\to T_1^\vee$.
Dualizing the restriction to $C$ gives
$\tau u^\vee=-g$, because $a^*\lambda=-\operatorname{id}_J$
\cite[Lemma 6.9 and Remark 6.10(c), p.~23]{MilneJV}.
Thus $-u^\vee:T_1\to T$ gives the required factorization.
For uniqueness, a homomorphism whose restriction is zero gives by
duality a map out of $\operatorname{Alb}(Y)$ vanishing on the generic
curve, hence on $Y$ by density and then on its Albanese variety.
This proves the trace assertion.  The finiteness of
the trace kernel is Conrad's Theorem 6.4(4); for the regular extension
$K/k$ its connectedness is Theorem 6.12 \cite{ConradTrace}.  These
statements allow the possible inseparable kernel in positive
characteristic, rather than assuming that the trace is a closed immersion.

By Poincar\'e complete reducibility, there is a homomorphism $s:J'\to J$
with $q s=[N]$ for some $N>0$.  Thus the cokernel of $J(K)\to J'(K)$ is
killed by $N$.  Also $T_K\to B$ is an isogeny, so an inverse isogeny up
to multiplication shows that $B(K)/\tau T(k)$ is torsion: here
$T(K)=T(k)$, since rational maps $\mathbf P^1\dashrightarrow T$ extend
and are constant.  Tensoring the resulting kernel and cokernel with
$\mathbf Q$ proves the displayed comparison.  In particular the argument
does not replace ``torsion cokernel'' by ``finite cokernel'' without a
finite-generation input.
The quotient $J'$ has zero trace by Conrad's Theorem 6.4, so his
Theorem 7.1 also shows that $J'(K)$ is finitely generated.  Thus the
torsion cokernel in this particular situation is indeed finite.
\end{proof}

\begin{theorem}[Proved: the rank formula, including reducible fibres]
\label{D:rank}
Let $r_v$ be the number of irreducible components of the fibre $Y_v$.
Then
\[
 \boxed{\rho(J')=\rho(Y)-2-\sum_v(r_v-1)
                 =\rho(X)+H^2-2-\sum_v(r_v-1).}
\]
Consequently the rank formula on cote 130, scan 14, holds under its
explicit irreducibility hypothesis.  We have inspected the scan: its
left-hand term is $\rho(\widetilde X)-2$, not $\rho(\widetilde X)-1$.
The integer entering the blow-up formula is $|Z|=H^2$, regardless of the
ambiguous gloss underneath the manuscript formula; it is not the number
of singular pencil members.
\end{theorem}
\begin{proof}
Every line bundle on $C$ is represented by a divisor, and its closure on
the regular surface $Y$ is Cartier.  Therefore
$\operatorname{Pic}(Y)\to\operatorname{Pic}(C)$ is surjective.
Its kernel consists of vertical divisor classes: if a divisor restricts
to a principal divisor on $C$, subtract the divisor of the same rational
function on $Y$.  The section identifies $\operatorname{Pic}^0(C)$ with
$J(K)$ and splits the degree map.  Hence the degree-zero subgroup of
$\operatorname{NS}(Y)$ maps onto $J(K)/\tau T(k)$, with kernel the
vertical classes.

The degree map is intersection with $F$ and is onto because $O\cdot F=1$,
so its kernel has rank $\rho(Y)-1$.  For each reducible fibre, omit the
component meeting $O$ and retain the other $r_v-1$ components.  Their
intersection matrix is negative definite; for a connected fibre this is
the usual negativity of the fibre intersection matrix after removal of
its one-dimensional radical.  Matrices at different fibres are mutually
orthogonal.  Together with $F$ these classes are independent: intersection
with $O$ first kills the coefficient of $F$, and negative definiteness
kills the other coefficients.  They span all vertical classes over
$\mathbf Q$, since the omitted component is obtained from the full fibre
and the other components.  The section forces its multiplicity to be
one.  Thus the vertical rank is $1+\sum_v(r_v-1)$.  Subtract this rank
and apply Lemma~\ref{D:fixed}.  Finally, blowing up $n$ distinct smooth
points adds $n$ independent exceptional classes to $\operatorname{NS}$.
\end{proof}

\subsection{The correction terms as inverse intersection matrices}

For a fibre $Y_v=\sum_{i=0}^{r_v-1}m_{v,i}\Theta_{v,i}$, number its
components so that $O$ meets $\Theta_{v,0}$, and put
\[
 M_v=(\Theta_{v,i}\cdot\Theta_{v,j})_{1\leq i,j<r_v},\qquad
 C_v=-M_v.
\]
An empty matrix contributes zero.  For a divisor $D$ on $Y$ with
$D\cdot F=0$, write
$b_v(D)=(D\cdot\Theta_{v,i})_{1\leq i<r_v}$ and set
\[
 \Phi_v(D)=-\sum_{i=1}^{r_v-1}(M_v^{-1}b_v(D))_i\Theta_{v,i},
 \qquad
 \nu_v(D,E)=b_v(D)^tC_v^{-1}b_v(E).
\]

\begin{proposition}[Proved, with the local-symbol comparison Known]
\label{D:local}
The corrected divisor $D+\Phi_v(D)$ has intersection zero with every
component of $Y_v$, and $\Phi_v(D)\cdot E=\nu_v(D,E)$.
For degree-zero divisors $D_K,E_K$ on $C$ with disjoint support, use their
closures to form these expressions.  The local N\'eron symbol, with
valuation of a uniformizer equal to one, is
\[
       (D_K,E_K)_v=(\overline D_K\cdot\overline E_K)_v
                          +\nu_v(\overline D_K,\overline E_K).
\]
If $e_v$ is the exponent of the component group of the N\'eron model of
$J$ at $v$, then its values belong to $e_v^{-1}\mathbf Z$.
Modulo $\mathbf Z$, the correction is the Grothendieck component-group
pairing, after identifying $J$ and $J^\vee$ by the positive canonical
principal polarization.  Thus the denominator assertion of cote 130,
scans 4--6, holds with the sharper exponent bound, and in particular
with $m_v=|\Phi_{J,v}|$.
\end{proposition}
\begin{proof}
The equations expressing zero intersection are exactly
$M_v a=-b_v(D)$, giving the formula for $\Phi_v(D)$.
They also imply zero intersection with the omitted component, since
$D\cdot F=0$ and $m_{v,0}=1$.  Multiplication by $b_v(E)^t$ gives the
assertion about $\nu_v$.  This calculation is valid with component
multiplicities, so it does not require a semistable fibre.

Complete the local ring of $\mathbf P^1$ at $v$.  The completed model
is proper, flat and regular over a complete discrete valuation ring,
has algebraically closed residue field, and has a section.  Completion
preserves the intersection numbers in question: lengths defining proper
intersections are unchanged by this faithfully flat completion.  Regularity
is preserved here because the model is of finite type over the excellent
base curve.  Thus every hypothesis of P\'epin's complete-DVR comparison is
satisfied, including smoothness and geometric connectedness of the generic
curve.  P\'epin's
Theorem 2.2.1 identifies the intersection construction with the N\'eron
symbol; its Corollary 2.2.2 gives the exponent bound
\cite[pp.~1323--1324]{Pepin2012}.  This is precisely the displayed
formula, since the vertical correction is the unique solution modulo
the full fibre of the equations just solved.

The component group is the torsion cokernel of the full fibre
intersection matrix.  Because $m_{v,0}=1$, it identifies with
$\mathbf Z^{r_v-1}/M_v\mathbf Z^{r_v-1}$.  The discriminant pairing of
this matrix is indeed the full-matrix pairing on degree-zero vectors:
if $m=(1,m_{v,1},\ldots,m_{v,r_v-1})^t$, the integral change of basis
$(e_0,\ldots,e_{r_v-1})\mapsto(m,e_1,\ldots,e_{r_v-1})$ has determinant
one and changes the full matrix by congruence to $0\oplus M_v$.
The first coordinate of a covector in the dual basis is its degree
$m^tb$, which vanishes for our divisors.  This gives both the asserted
torsion cokernel and the stated inverse-matrix formula, with no omitted
multiplicity factor.  The discriminant pairing of
this matrix sends the two component classes represented by $b,b'$ to
$b^tM_v^{-1}b'\bmod\mathbf Z$.  P\'epin's Theorem 4.1.1 identifies this
with Grothendieck's pairing using the \emph{opposite} of the principal
polarization; its index factor is one because of the section
\cite[pp.~1341--1343]{Pepin2012}.  Replacing that identification by the
positive polarization changes the sign and gives
$b^tC_v^{-1}b'=\nu_v\bmod\mathbf Z$.
The intrinsic pairing is the one defined as the Poincar\'e-biextension
extension obstruction in SGA 7, IX, 1.2.1 \cite{SGA7Pairing}.
\end{proof}

\subsection{The height formula and its normalization}

We use Shioda's bilinear normalization: if $\widehat h_\Theta$ is the
quadratic canonical height associated with a theta divisor, then
\[
 \langle P,Q\rangle_\Theta
  =\widehat h_\Theta(P+Q)-\widehat h_\Theta(P)-\widehat h_\Theta(Q),
 \qquad \langle P,P\rangle_\Theta=2\widehat h_\Theta(P).
\]
The negative sign in the intersection description below corresponds to
the positive principal polarization of a Jacobian.

\begin{theorem}[Proved: the height formula of cote 130]
\label{D:height}
For classes $P,Q\in J(K)$ represented by degree-zero divisors with
closures $D,E$ on $Y$, the height pairing is
\[
 \boxed{\langle P,Q\rangle_\Theta
             =-D\cdot E-\sum_v b_v(D)^tC_v^{-1}b_v(E).}
\]
It factors through $J(K)/\tau T(k)$ and, after tensoring with
$\mathbf Q$, through $J'(K)$.  If every fibre is irreducible, and
$\xi,\eta\in\operatorname{NS}(X)\otimes\mathbf Q$ satisfy
$\xi\cdot H=\eta\cdot H=0$, restriction and quotient define $i$ with
\[
                   \boxed{Q(i\xi,i\eta)=-\xi\cdot\eta.}
\]
Here $Q$ on the mobile quotient is the pairing transported from $J$ by
Lemma~\ref{D:fixed}, with the preceding bilinear normalization.
\end{theorem}
\begin{proof}
Put $\Phi(D)=\sum_v\Phi_v(D)$ and
\[
 D^\perp=D+\Phi(D)-((D+\Phi(D))\cdot O)F.
\]
It is orthogonal to $O$, to $F$ and to every vertical component.  The
last multiple of $F$ has no effect on pairings of degree-zero divisors.
Since components from different fibres do not intersect, expansion gives
\[
 D^\perp\cdot E^\perp
       =D\cdot E+\sum_v\nu_v(D,E).
\]
This also proves independence of the extensions: changing $D$ by a
vertical divisor has zero effect after projection, and changing it by
a principal divisor or an algebraically trivial divisor has zero
intersection with every divisor.  Closure of divisors and the kernel
calculation in Theorem~\ref{D:rank} show that these are exactly the
ambiguities that need to be removed.

The global N\'eron pairing is the sum of the local symbols; the product
formula kills principal-divisor changes.  We check the sign and the
normalization explicitly.  Lang's corollary to Theorem 5
\cite[p.~441]{LangHeights1964} identifies the global symbol of a
degree-zero cycle and an algebraically trivial line bundle $M$ on an
abelian variety with its canonical linear height, after changing his
$\log|\cdot|_v$ convention to our $\operatorname{ord}_v$ convention.
The scope over a field with a product-formula set of absolute values is
stated on p.~435; thus it applies to $K=k(t)$ with these valuations.
Let $a:C\to J$ be the Abel--Jacobi map defined by $O$.  For the positive
polarization $\lambda=\phi_{\mathcal O(\Theta)}$, autoduality gives
\[
             a^*\lambda=-\operatorname{id}_J.
\]
This is exactly Milne's Lemma 6.9 and Remark 6.10(c)
\cite[p.~23]{MilneJV}; the minus sign was checked on the rendered original
page, since text extraction omits it.  Consequently the line bundle on
$J$ pulling back to $\mathcal O_C(E)$ has class $-\lambda(Q)$.
Functoriality of the global symbol therefore makes its value on $D,E$
equal to $\widehat h_{-\lambda(Q)}(P)$.  Translation of the theta height
gives
\[
 \widehat h_{\lambda(Q)}(P)
  =\widehat h_\Theta(P+Q)-\widehat h_\Theta(P)-\widehat h_\Theta(Q).
\]
For example this follows by taking the linear part, as a function of
$P$, of $\widehat h_\Theta(P+Q)-\widehat h_\Theta(P)$; its constant part
is $\widehat h_\Theta(Q)$.  Thus the bilinear theta height is precisely
the negative of the sum of our local symbols.  Applying
Proposition~\ref{D:local} gives the boxed formula.
Equivalently it is $-D^\perp\cdot E^\perp$, the
intersection definition of the Mordell--Weil pairing.  For genus one,
this agrees term by term with Shioda's original Theorem 8.6,
equations (8.11)--(8.13), p.~228, which we inspected on the original
printed page \cite{Shioda1990}.  The same orthogonal-projection
construction for higher genus is stated in Shioda's Theorem 2.4
\cite[p.~248]{Shioda1992}; the local-symbol comparison above supplies its
canonical-height interpretation here.

The intersection form on the orthogonal complement of the span of
$O,F$ and the other fibre components is negative definite by the Hodge
index theorem.  Thus the displayed pairing is positive definite on the
quotient after removal of torsion.  It kills $\tau T(k)$ because
algebraically trivial divisors have zero intersection, and
Lemma~\ref{D:fixed} transports it to the mobile quotient.

Finally let $D=\pi^*\xi$ and $E=\pi^*\eta$.  Then $D\cdot F=E\cdot F=0$
and both have zero intersection with every exceptional section.
If all fibres are irreducible there are no matrices $C_v$, so the
correction terms vanish.  Birational invariance of intersection for
pullback gives $D\cdot E=\xi\cdot\eta$, proving the second box.
\end{proof}

\begin{remark}[Proved: the polarization on the quotient must be specified]
\label{D:quotient-polarization}
Let $\lambda:J\xrightarrow{\sim}J^\vee$ be the principal polarization.
The map
\[
       \mu=q\lambda^{-1}q^\vee:(J')^\vee\longrightarrow J'
\]
is the induced polarization on the dual quotient.  Its inverse
$\lambda'_{\mathbf Q}=\mu^{-1}$ is a positive rational polarization
$J'\to(J')^\vee$.  The form $Q$ in Theorem~\ref{D:height} is its height
pairing.  To check this, restrict $q$ to the $\lambda$-orthogonal
complement of $B$; it is an isogeny, and the pullback of
$\lambda'_{\mathbf Q}$ is the restriction of $\lambda$.  Functoriality
of canonical heights and the zero height of the trace give exactly the
transported form.  An integral polarization $N\lambda'_{\mathbf Q}$
gives $NQ$.  Consequently the manuscript's phrase ``canonical
polarization inherited from $J$'' needs this rational normalization:
an unspecified integral choice cannot justify an equality with coefficient
one.  There is a separate factor of two if one calls
$\widehat h_\Theta$ itself the bilinear height.
\end{remark}

\begin{proposition}[Proved counterexample when irreducibility is omitted]
\label{D:reducible-counterexample}
Even a Lefschetz pencil of positive-genus hyperplane sections need not
satisfy the uncorrected rank formula.  Over $\mathbf C$, take
\[
 E:y^2=x^3-x,\qquad X=E\times\mathbf P^1,
 \qquad H=\mathcal O_E(3O_E)\boxtimes\mathcal O_{\mathbf P^1}(1).
\]
Use homogeneous coordinates $[u:v]$ on the second factor and the pencil
\[
       (yu+xv)+t\bigl((x+2)u+(y+1)v\bigr)=0.
\]
It has six distinct transverse base points, six singular members each
with exactly one ordinary node, and every singular member has two
irreducible components.  Its generic Jacobian is the constant elliptic
curve $E_K$, so $J'=0$.  Nevertheless
$\rho(X)+H^2-2=2+6-2=6$.
\end{proposition}
\begin{proof}
The three functions $1,x,y$ give the complete sections of
$\mathcal O_E(3O_E)$, which is very ample.  Thus $H$ is very ample and
$H^2=6$.  At a point of $E$ where the two coefficients of a pencil member
do not both vanish, that member is the graph of a unique point of
$\mathbf P^1$.  A vertical component occurs precisely when
\[
 y+t(x+2)=0,\qquad x+t(y+1)=0.
\]
There are no solutions for $t=1,-1$; otherwise these equations give
\[
 x=\frac{t(2t-1)}{1-t^2},\qquad
 y=\frac{t(t-2)}{1-t^2}.
\]
Substitution into $y^2=x^3-x$ gives
\[
        p(t)=-t(7t^5-15t^4+13t^3+t^2-6t+1)=0.
\]
Its discriminant is $-101056167\ne0$, and
$p(1)=-1$, $p(-1)=-27$.  These exact identities are checked in
\texttt{code/D\_pencil\_checks.py}, with full output retained.
There is no solution at $t=\infty$, since $x=-2,y=-1$ is not on $E$.
There is no base point or vertical component above $O_E$: in the local
trivialization of $\mathcal O_E(3O_E)$, the coefficient matrix there is
diagonal with nonzero entries.

The displayed equations solve uniquely for $x,y$ at each root of $p$,
so they give six distinct points of $E$, with distinct pencil parameters.
The coefficient matrix
$\left(\begin{smallmatrix}y&x\\x+2&y+1\end{smallmatrix}\right)$ has rank
one at each such point and never has rank zero.  The simple-root
calculation says that its determinant has a simple zero on $E$.
Consequently its kernel gives one base point above each point of $E$,
and the two generators have independent differentials there.  This
verifies transversality of the six base points.

At each exceptional parameter the member consists of one vertical
$\mathbf P^1$ and the residual graph of a morphism $E\to\mathbf P^1$.
In a local coordinate $s$ on $E$ at the common zero, its equation is
$s(a(s)u+b(s)v)=0$, with $(a(0),b(0))\ne(0,0)$.  The two components
therefore meet transversely at one point, giving one ordinary node.
The determinant transversality also places the pencil base point on the
vertical component away from this node.  All other members are graphs
and are smooth copies of $E$.  Blowing up the six transverse base points
gives the required Lefschetz fibration, with six two-component fibres.
Finally $\operatorname{Pic}(E\times\mathbf P^1)=
\operatorname{Pic}(E)\oplus\mathbf Z$, so $\rho(X)=2$; the generic curve
is $E_K$, whence its mobile Jacobian is zero.  The corrected formula
gives $2+6-2-6=0$, as required.
\end{proof}

The same script also checks the smaller genus-zero example
$x(y-z)-t\,y(x-z)=0$ on $\mathbf P^2$ with $H=2L$:
four transverse base points, three reducible nodal conics, and rank
$1+4-2-3=0$.  Its uncorrected value would be $3$.

\begin{remark}[Proved: the familiar elliptic formula]
\label{D:elliptic}
For a relatively minimal elliptic surface and sections $P,Q$, take
$D=(P)-(O)$, $E=(Q)-(O)$.  Since $(O)^2=-\chi(\mathcal O_Y)$ and the
nonidentity components do not meet $O$, Theorem~\ref{D:height} becomes
\[
 \langle P,Q\rangle=\chi(\mathcal O_Y)+(P\cdot O)+(Q\cdot O)
       -(P\cdot Q)-\sum_v\operatorname{contr}_v(P,Q),
\quad
 \operatorname{contr}_v(P,Q)=b_v(P)^tC_v^{-1}b_v(Q).
\]
This is exactly Shioda's (8.11)--(8.13), including the sign of the local
correction.  The rank calculation and the general projection formula
above do not require relative minimality; that condition is only used
here for the standard self-intersection identity.
\end{remark}
